\documentclass[10pt, aspectratio=169]{beamer}

\usetheme{Madrid}
\usecolortheme{dolphin}

\usepackage{ctex}
\usepackage{xeCJK}
\usepackage{amsmath, amssymb}
\usepackage{graphicx}
\usepackage{listings}
\usepackage{tikz}
\usepackage{xcolor}

\title{408考研零基础 --- 数据结构}
\subtitle{1-6 栈}
\author{}
\date{}

\setbeamertemplate{page number in head/foot}{}

\begin{document}


\begin{frame}{本节目标}
    \begin{enumerate}
        \item 理解栈的定义、逻辑特征及其 ADT 定义
        \item 掌握顺序栈和链栈的存储结构与基本操作实现
        \item 理解栈在括号匹配、表达式求值等场景中的应用原理
    \end{enumerate}
\end{frame}

\begin{frame}{栈的定义与 ADT}
    \textbf{栈}：只允许在一端进行插入和删除操作的线性表
    \begin{itemize}
        \item 栈顶 --- 允许操作的一端
        \item 栈底 --- 不允许操作的一端
        \item 原则：后进先出（LIFO）
    \end{itemize}
    \vspace{0.5em}
    \textbf{基本操作}
    \begin{itemize}
        \item InitStack(\&S) --- 初始化
        \item Push(\&S, e) --- 入栈
        \item Pop(\&S, \&e) --- 出栈
        \item GetTop(S, \&e) --- 取栈顶
        \item StackEmpty(S) --- 判栈空
    \end{itemize}
\end{frame}

\begin{frame}[fragile]{顺序栈}
    \begin{lstlisting}[language=C, basicstyle=\ttfamily\small]
#define MaxSize 50
typedef struct {
    ElemType data[MaxSize];
    int top;
} SqStack;
    \end{lstlisting}
    \vspace{0.5em}
    \textbf{核心操作}
    \begin{itemize}
        \item 初始化：$S \to \text{top} = -1$
        \item 入栈：$S \to \text{data}[++S \to \text{top}] = e$
        \item 出栈：$e = S \to \text{data}[S \to \text{top}--]$
        \item 判空：$S \to \text{top} == -1$
    \end{itemize}
    \vspace{0.3em}
    入栈出栈时间复杂度 $O(1)$
\end{frame}

\begin{frame}{共享栈}
    \textbf{两个栈共享同一数组空间}
    \vspace{0.5em}
    \begin{itemize}
        \item 左栈栈底在 0 端，$S_1 \to \text{top} = -1$
        \item 右栈栈底在 MaxSize-1 端，$S_2 \to \text{top} = \text{MaxSize}$
        \item 判满条件：$S_1 \to \text{top} + 1 == S_2 \to \text{top}$
    \end{itemize}
    \vspace{0.5em}
    优点：更有效地利用存储空间
\end{frame}

\begin{frame}{链栈}
    \textbf{操作受限的单链表}
    \vspace{0.5em}
    \begin{itemize}
        \item 入栈：在头结点之后插入
        \item 出栈：删除头结点之后的结点
        \item 无需预分配空间
        \item 不会出现栈满
        \item 时间复杂度 $O(1)$
    \end{itemize}
    \vspace{0.5em}
    \textbf{入栈}
    \[
    s \to \text{next} = S \to \text{next};\quad S \to \text{next} = s;
    \]
    \textbf{出栈}
    \[
    p = S \to \text{next};\quad S \to \text{next} = p \to \text{next};\quad \text{free}(p);
    \]
\end{frame}

\begin{frame}{栈的应用 --- 括号匹配}
    \textbf{算法步骤}
    \begin{enumerate}
        \item 遍历表达式
        \item 遇到左括号 $\to$ 入栈
        \item 遇到右括号 $\to$ 检查栈顶是否为对应左括号
        \begin{itemize}
            \item 匹配 $\to$ 出栈
            \item 不匹配 $\to$ 失败
        \end{itemize}
        \item 遍历结束，栈为空则匹配成功
    \end{enumerate}
\end{frame}

\begin{frame}{栈的应用 --- 表达式求值}
    \textbf{中缀转后缀}
    \begin{itemize}
        \item 操作数 $\to$ 直接输出
        \item 运算符 $\to$ 与栈顶比较优先级
        \begin{itemize}
            \item 优先级高 $\to$ 入栈
            \item 优先级低/相等 $\to$ 出栈并输出
        \end{itemize}
    \end{itemize}
    \vspace{0.5em}
    \textbf{后缀计算}
    \begin{itemize}
        \item 操作数 $\to$ 入栈
        \item 运算符 $\to$ 弹出两个操作数计算，结果入栈
    \end{itemize}
\end{frame}

\begin{frame}{例题 --- 十进制转八进制}
    \textbf{除基取余法}
    \vspace{0.5em}
    十进制 1348 转八进制：
    \begin{itemize}
        \item $1348 \div 8 = 168 \cdots\cdots 4$
        \item $168 \div 8 = 21 \cdots\cdots 0$
        \item $21 \div 8 = 2 \cdots\cdots 5$
        \item $2 \div 8 = 0 \cdots\cdots 2$
    \end{itemize}
    \vspace{0.5em}
    余数入栈 $\to$ 出栈得 2504
    \[
    2 \times 8^3 + 5 \times 8^2 + 0 \times 8^1 + 4 \times 8^0 = 1348
    \]
\end{frame}

\begin{frame}{本节总结}
    \begin{enumerate}
        \item 栈是 LIFO 的受限线性表
        \item 顺序栈：数组实现，$O(1)$，需预分配空间
        \item 链栈：链表实现，$O(1)$，不会栈满
        \item 应用：括号匹配、表达式求值、函数递归
    \end{enumerate}
\end{frame}

\end{document}
